%% degrees-and-priority-continuation.tex
%% Revised proofs for the editorial degree and priority continuation.

% Necessity: the later finite-injury result retains proof detail not
% present in the exact inherited source or its reader-facing proofs.
% Definition: this annex supplies complete proofs for jump extensivity,
% Kleene--Post, and Friedberg--Muchnik, retaining the historical strategy core
% where it remains correct.
% Standing: the annex continuation is administrative navigation syntax, not a
% new proof status, theorem family, mathematical relation, or admission claim.
% Nonredundancy: exact inherited proofs remain in the body; this file supports
% only the distinct later degree and priority results.
% Use: annex/p04 routes the file once and the body continuation points to its
% retained priority-construction proof item.
\noindent\textit{Reader note.}
These complete proofs support only the separately authored degree and priority
continuation; they do not replace any inherited source proof.

\begin{annexitem}{annex:p4.degrees.1}{thm:p04-jump-extensive}{Every oracle reduces to its jump}
\begin{proof}
Given \(n\), effectively produce the code \(q(n)\) of an oracle program that
ignores its numerical input, asks whether \(n\) belongs to the oracle, writes
\(0\) and accepts when the answer is yes, and otherwise loops forever.  The
fixed row encoding makes this compiler primitive recursive: insert the binary
numeral for \(n\) into one fixed query-program template, serialize its finite
rows, and place the result in the designated program-code slot.  Thus \(q\) is
total computable without invoking the sketched relativized parameter theorem.
Because the program ignores its self-input,
\[
 n\in A
 \quad\Longleftrightarrow\quad
 \varphi^{A}_{q(n)}(q(n))\mathord\downarrow
 \quad\Longleftrightarrow\quad q(n)\in A'.
\]
Thus \(q\) is a many-one reduction \(A\leq_mA'\).
\end{proof}
\end{annexitem}

\begin{annexitem}{annex:p4.degrees.2}{thm:p04-kleene-post}{Kleene--Post incomparability theorem}
\begin{proof}
Use the fixed oracle-program numbering \((\varphi_e^B)_{e<\omega}\).  Build
increasing strings \(\sigma_s,\tau_s\in2^{<\omega}\); their unions are the
characteristic functions of \(A\) and \(B\).  Meet, alternately,
\[
 R_e:\chi_A\not\simeq\varphi_e^B,
 \qquad S_e:\chi_B\not\simeq\varphi_e^A.
\]
Every candidate Turing reduction occurs in this list.  Given an oracle decider,
compile it to a numerical oracle program that writes \(1\) and accepts after
the simulated accept, and writes \(0\) and accepts after the simulated reject.
Its computed function is the target characteristic function.  The wrapper is
a primitive-recursive finite-table transformation, and the fixed serialization
of row tags, states, symbols, targets, and query rows produces its valid code.
For a finite oracle string \(\tau\), the notation
\(\varphi_e^\tau(x)\mathord\downarrow=i\) means that the finite computation
halts with value \(i\) and every oracle query is below \(\operatorname{len}(\tau)\),
answered by \(\tau\).

At the \(R_e\)-stage choose
\(x\geq\operatorname{len}(\sigma_s)\), pad \(\sigma_s\) with zeros below
\(x\), and ask \(\varnothing'\) whether
\[
 \exists\tau'\supseteq\tau_s\;
       \varphi_e^{\tau'}(x)\mathord\downarrow=0.
\]
This is a c.e. question about finite strings and finite computations.  By the
cited completeness theorem for \(K\),
Theorem~\ref{thm:arithmetical-hierarchy-classification}, it reduces many-one
to \(K\).  Given an
ordinary index \(p\), a primitive-recursive table compiler produces a
query-free oracle program that ignores its input and simulates
\(\varphi_p(p)\); its code belongs to \(\varnothing'\) exactly when
\(p\in K\).  Hence the question is decidable by \(\varnothing'\), without
using the sketched convention bridge.  If the
answer is yes, effectively find the least witnessing \(\tau'\), adopt it, and
set \(\sigma_{s+1}(x)=1\).  If the answer is no, retain \(\tau_s\) and set
\(\sigma_{s+1}(x)=0\).  Pad both new strings so their lengths exceed \(x\).
The \(S_e\)-stage interchanges \(\sigma\) and \(\tau\).  Also pad at every
stage beyond the stage number, so both unions are total.

Each requirement, once treated, is permanent.  In the yes case the chosen
finite oracle computation remains convergent with value \(0\) under every
later oracle extension, while the target bit is \(1\).  In the no case no
later extension can make that value \(0\); the final computation therefore
diverges or has a value different from the target bit \(0\).  Hence all
\(R_e,S_e\) hold, so \(A\not\leq_TB\) and \(B\not\leq_TA\).

The entire stage construction is computable in \(\varnothing'\).  To decide
the \(n\)-th bit of either limit, run it until the corresponding string has
length greater than \(n\); thus \(A,B\leq_T\varnothing'\).  Neither is
computable, since a computable set reduces to every oracle.  Neither has
degree \(\mathbf0'\), because if, for example,
\(\varnothing'\leq_TA\), then \(B\leq_T\varnothing'\leq_TA\), contrary to
incomparability.  Their degrees are therefore incomparable and strictly
between \(\mathbf0\) and \(\mathbf0'\).
\end{proof}
\end{annexitem}

\begin{annexitem}{annex:p4.degrees.4}{thm:p04-friedberg-muchnik}{Friedberg--Muchnik theorem}
\begin{proof}
List
\[
 R_e:\varphi_e^B\not\simeq\chi_A,
 \qquad S_e:\varphi_e^A\not\simeq\chi_B
\]
in the priority order \(R_0>S_0>R_1>S_1>\cdots\).  Use the fixed effective
oracle-program numbering.  Write
\(\varphi^{B_s}_{e,s}(x)\mathord\downarrow=y\) when the first \(s\) steps of
program \(e\), with each oracle query answered by membership in the finite set
\(B_s\), halt accepting with valid output \(y\).  Its use is one plus the
largest queried number, or \(0\) if no query occurs.  A finite computation is
determined by the oracle below its use: if all later approximations agree with
\(B_s\) there, it is the final \(B\)-oracle computation; conversely every
halting final computation appears at a sufficiently late stage once its finite
oracle segment stabilizes.  These claims follow by induction through the
finitely many transition steps.

Every candidate Turing reduction appears among the requirements.  Wrap its
oracle decider so that a simulated accept writes \(1\) and accepts, while a
simulated reject writes \(0\) and accepts; the resulting total numerical
function is the target characteristic function.  Normalizing and serializing
that finite table is primitive recursive and does not rely on the sketched
relativized-adequacy theorem.

At stage \(s\), a finite state records
\(A_s,B_s\), and, for each of the first \(s\) strategies, either no witness or
one witness, a satisfaction flag, and a restraint bound; bound \(0\) encodes
the empty restraint.  Visit the first \(s\) strategies in priority order.
This is a fair computable scheduler because every strategy is visited at every
sufficiently large stage.

An initialized \(R_e\) chooses a fresh witness \(x>s\), larger than every
active higher-priority restraint and every reserved witness, and keeps
\(x\notin A_s\).  Search the finite stage computation for
\(\varphi_{e,s}^{B_s}(x)\mathord\downarrow=0\).  When it appears, let \(u\)
be one plus the largest oracle query, or \(0\) if there was no query; enumerate
\(x\) into \(A\), impose the \(B\)-restraint \(u\), mark \(R_e\) satisfied,
and initialize every lower-priority strategy.  Initialization cancels its
witness, flag, and restraint but never removes an enumerated number.  A lower
strategy chooses all later witnesses above the active higher restraints, so
it cannot change the protected oracle segment.  The \(S_e\)-strategy is the
same construction with \(A,B\) interchanged.  Each stage performs finitely
many bounded simulations and finite-state updates, so the two stage unions are
c.e.; the fixed scheduler and universal numbering give their enumeration
indices effectively.

We verify requirements by induction through the priority order.  Assume all
higher strategies have stabilized.  They act only finitely often, so the
current strategy is initialized only finitely often and eventually chooses a
final witness.  Suppose for contradiction that
\(\varphi_e^B=\chi_A\).  If \(R_e\) acts on its final witness, the preserved
computation has value \(0\) while that witness enters \(A\), a contradiction.
If it never acts, the final witness stays outside \(A\), so total equality
would give a convergent value \(0\).  That computation has finite use.  Since
witnesses are chosen above their stages, only finitely many later lower-priority
enumerations can occur below that use; after the oracle segment stabilizes,
the stage simulation sees the value \(0\) and \(R_e\) acts, again a
contradiction.  Thus \(R_e\) is met; the proof for \(S_e\) is symmetric.

Every oracle program occurs in the fixed enumeration, so the requirements
give \(A\not\leq_TB\) and \(B\not\leq_TA\).  Neither set is computable.  Every
c.e. set is Turing reducible to \(\varnothing'\): from an enumeration index
and a number \(n\), compile a query-free oracle program that writes \(0\) and
accepts exactly when that enumeration lists \(n\), and ask
\(\varnothing'\) whether its code is in
the jump.  If either constructed set had
degree \(\mathbf0'\), the other would reduce to it, contradicting
incomparability.  Hence both c.e. degrees lie strictly between
\(\mathbf0\) and \(\mathbf0'\).
\end{proof}
\end{annexitem}
